\chapter{从基因到性状，中间隔着整个生命}

\begin{kaichang}
打开购物网站，搜“基因检测”，你会看到这样一门生意：寄去一管唾液，几周后收到一份几十页的报告——你的祖源成分、你喝酒上不上脸、你的咖啡因代谢快不快、你是“爆发力型”还是“耐力型”运动员、你将来患某些疾病的风险高不高。广告词写得斩钉截铁：基因里藏着答案。

现在先别下单，做一道思考题。假设这份报告完全诚实、技术也完全可靠，它能把你的 DNA 序列读得一字不差。那么，它能告诉你的东西，和不能告诉你的东西，分界线画在哪里？它能预言你会成为一个什么样的人吗？

先猜一个答案。这一章结束时，你会有办法给那份报告里的每一句话打分——而且你会发现，报告里最误导人的，往往不是测错的地方，而是那种“基因即命运”的说话方式。
\end{kaichang}

\section{性状的两种脾气：档位和旋钮}

先做现象层该做的事：不看基因，只看人。

看看你身边的人。有些性状像电灯开关，只有几个\textbf{档位}：ABO 血型，你要么是 A 型、B 型、AB 型、O 型，没有“A 型半”；你吃蚕豆会不会出事（蚕豆病）、你的耳朵耳垢是干是湿，也基本是离散的几种。另一些性状则像调光旋钮，是\textbf{连续}的：身高从一米四到两米，几乎任何数值都有人站在那儿；肤色从很深到很浅，是一道没有断口的光谱；血压、体重、跑步成绩，全都如此。

再看一个更微妙的现象。同卵双胞胎来自同一个受精卵，DNA 序列几乎一模一样。他们确实像——站在一起常常让人分不清。但只要熟识他们的人都知道：两个人性格有别、发型有别，连指纹都不完全相同；如果一个后来学了钢琴、一个去踢了足球，十年之后连体态都看得出差别。序列相同，成品却不是一个模子里倒出来的。

还有时间维度。你的曾祖父辈的平均身高，比今天的同龄人矮一截——这不是因为他们的基因在一个世纪里换了（基因频率变化没那么快），而是因为营养、医疗、童年疾病负担变了。同一份“图纸”，在不同的施工条件下，盖出了不同高度的楼。

这三组现象——档位与旋钮、双胞胎的差异、隔代的身高变化——是本章要解释的全部事实。而解释的钥匙，藏在从 DNA 到性状那条长得超乎想象的链条里。

\section{一个字母怎样变成一场病}

把镜头推进到粒子层。为了让链条的每一环都看得见，我们全程跟踪一个真实案例：镰状细胞贫血。它是人类理解得最透彻的遗传病之一，也是本章的“导游”。

故事从血液开始。你的红细胞里没有细胞核，几乎被一种蛋白质塞满——血红蛋白（第 48 章见过它）。一个血红蛋白分子由四条链组成，其中两条叫 $\beta$ 珠蛋白链。$\beta$ 珠蛋白的“配方”写在 DNA 上：经过转录（第 68 章）和翻译（第 69 章），核糖体按密码子把氨基酸一个个串起来。配方的第 6 个位置，正常情况下是一个叫谷氨酸的氨基酸，它的侧链带负电、亲水，安静地待在蛋白质表面。

在有些人的 DNA 里，这个位置对应的那个碱基，从 A 变成了 T——第 72 章讲过的碱基替换，三十亿个字母里的一次。就这一次。密码子从 GAG 变成 GUG，翻译机器读到 GUG，放进去的就不是谷氨酸，而是缬氨酸。缬氨酸的侧链疏水——怕水，像一小滴油。

一个“油点”出现在蛋白质表面，会有什么后果？第 48 章讲过疏水作用：怕水的部分在水环境里倾向于躲开、抱团。单个缬氨酸的替换平时掀不起大浪，但当红细胞把氧气卸下、血红蛋白改变形状时，这个疏水点恰好暴露出来，于是一个又一个血红蛋白分子“油点贴油点”，首尾相接，聚合成坚硬的长纤维。纤维在红细胞内部越攒越长，把原本圆饼状、柔软的红细胞，硬生生撑成弯曲的镰刀形。

链条还没走完，现在轮到细胞层和组织层。镰刀形的红细胞又硬又脆：它们在狭窄的毛细血管里会卡住、堆积，堵住血流；它们的寿命从大约 120 天缩短到十几天，骨髓造新的来不及补上——于是贫血、组织缺氧、阵发性剧痛，长年累月还会损伤脾、肾等器官。一个碱基的差异，就这样一级一级放大，变成一个人的病痛：

\begin{center}
\begin{tikzpicture}[scale=1]
\node[align=center] at (0,0) {\small 一个碱基 A$\rightarrow$T};
\draw[->, thick] (1.7,0) -- (2.3,0);
\node[align=center] at (3.9,0) {\small 一个氨基酸被换掉};
\draw[->, thick] (5.5,0) -- (6.1,0);
\node[align=center] at (7.6,0) {\small 蛋白质聚成纤维};
\draw[->, thick] (9.1,0) -- (9.7,0);
\node[align=center] at (11.1,0) {\small 红细胞变形};
\draw[->, thick] (5.5,-0.7) -- (5.5,-1.3);
\node[align=center] at (5.5,-1.9) {\small 血管堵塞、贫血、疼痛、器官损伤};
\node[align=center] at (5.5,-2.7) {\small 示意图：每一步都是真实的物理化学过程，不是比喻};
\end{tikzpicture}
\end{center}

请注意这条链条的每一环都不是“信息”在直接起作用，而是货真价实的物质和能量：侧链的亲水与疏水（第 23 章的水、第 47 章的脂质逻辑）、蛋白质的形状变化、细胞的力学性质、血流的几何。基因只写了第一行，剩下的剧情全部由物理和化学接着演。

\section{基因型到表型，隔着一张网}

有了镰状细胞贫血这个完整样本，我们可以把规则层说清楚了。规则只有一句话：\textbf{基因型是输入，表型是输出，但中间的网络由蛋白质、细胞、组织、环境和历史共同组成，谁也无法跳过。}

先看“网络”意味着什么。从 DNA 到你的身高，中间至少有这些层：基因是否表达、表达多少（第 70 章的调控）；蛋白质造出来后怎样折叠、修饰、和谁搭档；蛋白质组成的细胞怎样分裂、迁移、死亡（第 75 章会专门讲发育）；细胞组成的组织怎样和激素、营养、运动相互作用。链条上每一层都有别的基因插手，也都有环境插手。

这就引出三组重要概念。

第一组：\textbf{单基因与多基因}。像 ABO 血型、镰状细胞贫血这样主要由一个基因座决定的性状，是少数派。身高这样的数量性状是多数派：全基因组研究已经找到成百上千个与身高相关的基因座，每个的贡献都只有几毫米甚至更少，像一支几百人的合唱队，没有人独唱。肤色也一样，由几十个基因共同调配黑色素的种类和产量——所以孩子的肤色常常是父母的“混音”，而不是非此即彼的选择题。

第二组：\textbf{外显率与表现度}。两个人带着一模一样的致病基因型，结局可能不同。外显率说的是：带这个基因型的人里，有多大比例真的会表现出相应性状。有些基因型外显率是百分之百，有些不是——比如 BRCA1 的某些致病突变显著提高乳腺癌风险，但携带者中终生患病的比例大约是六成到七成，不是十成。表现度说的是：即使都表现出来，程度也可以天差地别。同样是镰状细胞贫血的纯合患者，有人幼年即重病，有人症状较轻——差别来自别的基因、胎儿血红蛋白的残留量、感染史、医疗条件。基因型相同，剧本相同，演出效果仍看剧场和天气。

第三组：\textbf{基因互作}。一个基因的效果，常常取决于别的基因的版本。教科书里常被当作“单基因性状”的卷舌能力（能不能把舌头卷成筒），其实是反例的好教材：双生子研究显示，同卵双胞胎中一个会卷舌、一个不会的情况并不少见，练习也能改变——它根本不是一道单基因判断题。你体内没有哪个基因是独自发言的，它们更像在一张谈判桌上互相投票。

最后是环境，而且环境不只是“外部”。营养、童年感染、睡眠、运动、长期压力（压力通过激素一路传导到细胞核门口，用的正是第 58 章的信号逻辑）、母亲孕期的营养状况、甚至你的发育顺序（先长了什么、后长了什么），都会在链条的每一层留下痕迹。这些痕迹有些还能以“序列不变、使用方式改变”的方式被细胞记住——那叫表观遗传，第 74 章专门讲它。

把上面的规则合起来用一次，看一个你也许没想到的案例：喝牛奶。几乎所有人类婴儿都能消化乳汁里的乳糖（第 46 章讲过乳糖不耐受），因为肠道里有乳糖酶。但断奶之后，大多数人的乳糖酶基因会“调低音量”（第 70 章的调控）——成年人喝奶就腹胀、腹泻，这才是人类的“默认设置”。而在有数千年养牛饮奶传统的族群（北欧、东非草原等地）里，一种调控区的变异让这盏灯终身不灭，携带者能从牛奶里获得宝贵的热量、钙和水分。乳糖耐受这个性状，是基因调控变异与饮食文化长期共舞的产物：序列负责硬件，调控负责开关，文化负责环境——缺了任何一环，你都解释不了为什么你同桌喝奶没事、你喝奶跑厕所。

\begin{qiaoliang}
有一个词你必须学会正确使用：\textbf{遗传率}。你一定见过这样的说法——“身高的遗传率是 $80\%$”。大多数人第一反应是：“我的身高有八成长在基因里。”这是错的，而且是本章最想纠正的错误。

遗传率是一个\textbf{关于群体}的统计量：在一个特定群体、特定环境里，人与人之间性状的\textbf{差异}，有多大比例与基因差异相联系。$80\%$ 的意思是“这群人身高的参差不齐，八成能用他们的基因差异来解释”——而不是“你个人的身高八成由基因决定”。个人没有“差异的分解”这回事，你不是一群人。

更妙的是，遗传率会随环境改变。设想两块田：第一块田土壤均匀、灌溉一致，种上不同品种的玉米——这时高矮差异几乎全由品种（基因）决定，遗传率接近 $100\%$。第二块田种同一品种，但一半施肥一半不施——这时高矮差异几乎全由环境决定，遗传率接近 $0$。同一批种子，换一个环境布局，遗传率就变了。这就是为什么“身高遗传率 $80\%$”和“过去一百年平均身高猛涨”完全不矛盾：增长来自整体环境改善（营养、医疗），而同一时代人与人之间的参差仍主要由基因差异解释。一个统计量，两回事。
\end{qiaoliang}

\section{把链条压缩成符号}

现在轮到符号登场。第 63 章你学过用字母代表等位基因、用棋盘格算概率，这套符号在本章的导游案例上用起来格外顺手。

$\beta$ 珠蛋白基因座有两个常见等位基因：正常的记作 $\mathrm{Hb^{A}}$，镰状的记作 $\mathrm{Hb^{S}}$。每个人都有两份拷贝（分别来自父母），于是有三种基因型、三种命运：

\begin{table}[h]
\centering
\begin{tabular}{@{}lll@{}}
\toprule
基因型 & 蛋白质层面 & 个体层面（表型） \\
\midrule
$\mathrm{Hb^{A}Hb^{A}}$ & 血红蛋白全部正常 & 红细胞正常，不贫血 \\
$\mathrm{Hb^{A}Hb^{S}}$ & 两种血红蛋白混合 & 通常健康，红细胞基本不变形 \\
$\mathrm{Hb^{S}Hb^{S}}$ & 血红蛋白会聚合 & 镰状细胞贫血 \\
\bottomrule
\end{tabular}
\end{table}

符号的压缩力惊人：一个 $\mathrm{Hb^{S}}$ 就装下了“碱基 A 变 T、谷氨酸变缬氨酸、蛋白质聚合”这整条链。但压缩是要付出代价的——表格里那一格“通常健康”，把表现度、环境、基因互作全部藏进了省略号。两个 $\mathrm{Hb^{A}Hb^{S}}$ 的人结婚，棋盘格（第 63 章）告诉你每个孩子有 $1/4$ 概率是 $\mathrm{Hb^{S}Hb^{S}}$——这是真实的概率，可它只说了概率，没有说那个孩子的症状会多重、生活质量会怎样。符号给出的是群体的赌注，不是个人的判决书。

顺便做个数量级估算，感受一下“符号压缩了多少信息”。两个毫无血缘关系的人，DNA 序列大约有 $99.9\%$ 相同。人类基因组约 $3\times 10^{9}$ 个碱基，也就是说任意两人之间大约有
\[3\times 10^{9}\times 0.1\% = 3\times 10^{6}\]
即约三百万个位点不同。这三百万处差异里，绝大多数落在不编码、不调控的区域，或者干脆不产生可察觉的影响；真正让你和你同桌不一样的，是其中一小撮差异与整个环境的合奏。下次听到“找到某某性状的基因”时，记得这个背景：那是从三百万处差异里捞针，而且捞到的针几乎从不是单独作案。

\section{科学家怎样把一个症状追到一个碱基}

\begin{zhengju}
这条从碱基到病痛的链条，不是谁坐在书房里推出来的，它是一环一环被实验拧紧的，前后花了将近半个世纪。

第一环在显微镜下。1910 年，芝加哥医生詹姆斯·赫里克收治了一位来自加勒比海的年轻患者，他有严重的贫血和黄疸。赫里克把他的血做成涂片，在显微镜下看到了从没见过的东西：许多红细胞不是圆饼形，而是“细长的镰刀形、新月形”。他把这种形态记录下来，命名为镰状细胞贫血。注意，这时他能说的只有“细胞形状异常”——基因、DNA、蛋白质，一个都还没登场。

第二环等了三十九年。1949 年，莱纳斯·鲍林的团队用电泳做实验：把血红蛋白放在电场里，看它迁移的快慢——迁移率取决于分子表面带的电荷。结果干净利落：患者的血红蛋白和正常人的血红蛋白迁移率不同，杂合子（$\mathrm{Hb^{A}Hb^{S}}$）的血样里两种都有。仅仅一个电荷差异，就改变了整个分子在电场中的行为。鲍林把镰状细胞贫血称为第一种“分子病”——这个词本身就是一次思想革命：一种复杂的人类疾病，可以被定位到一个分子的物理性质上。

第三环更精细。1956 年，弗农·英格拉姆用“指纹图谱法”把正常和镰状血红蛋白分别切成小段、铺在纸上双向分离，发现所有片段几乎完全重合，只有一个小小的肽段位置不同。再分析这个肽段，结论出来了：整条链约 150 个氨基酸里，只有第 6 位从谷氨酸换成了缬氨酸。一个氨基酸，解释了电场里的迁移差异；今天我们知道，它背后是 DNA 上的一个碱基。

还有第四环，它不在实验室里，在地图上。有人注意到：镰状等位基因的高发地区，和疟疾肆虐的地区——撒哈拉以南非洲、地中海沿岸、印度部分地区——惊人地重合。1954 年，安东尼·艾利森提出并验证了大胆的猜想：携带一个 $\mathrm{Hb^{S}}$ 拷贝的杂合子，对重症疟疾有更强的抵抗力。地图上的重合不是巧合，是自然选择留下的指纹。这个环先按下不表——它是本章末尾伏笔的全部内容。

回头看这个故事的方法：形态观察提出“哪里不对”，物理测量锁定“分子带电不同”，化学分析精确到“一个氨基酸”，群体调查追问“为什么它这么常见”。每一环的答案都制造了新问题，而替代解释——比如“是感染导致的”“是营养不良”——都被对照实验一一排除。没有人一步登天，但每一环都拧得实实在在。
\end{zhengju}

\section{这套模型什么时候会失灵}

“基因 $\rightarrow$ 蛋白质 $\rightarrow$ 细胞 $\rightarrow$ 个体”这条链是真实存在的，但它有明确的适用边界，越过边界就会出错。

边界一：\textbf{单基因不等于单结局}。即使是镰状细胞贫血这样“最干净”的遗传病，外显率接近百分之百，表现度仍然因人而异。用它做基因型层面的预测很可靠，用它预言一个人具体的病程，就不那么可靠了。

边界二：\textbf{多基因性状只能给概率，而且概率会过期}。对身高、血压、糖尿病风险这类性状，科学家把几百上千个基因座的小效应加总成“多基因风险评分”。这类评分在群体层面有用（比如找出风险最高的一小撮人重点筛查），但对个人，它的误差大得惊人——因为每个基因效应都小、基因间互作被简化、环境权重巨大。更要命的是，这些评分大多基于特定祖源人群的数据算出，换一个人群就明显失准。拿它给个人“算命”，是把群体统计错当成个人剧本。

边界三：\textbf{相关不是命运，概率不是期限}。“风险提高三倍”听着吓人，但如果原来的风险是十万分之一，三倍也只是十万分之三。剂量、基线、证据等级——第二册里反复强调的这三样东西，在基因报告上一样都不能少。

\begin{wuqu}
“既然身高的遗传率高达 $80\%$，说明身高基本由基因决定，营养和干预没多大用。”——这是把群体统计偷换成个人因果的经典错误。回忆那两块田：同一块田里，品种差异决定高矮（遗传率高），但只要给整块田换上更好的水土，所有玉米都会长高。现实版证据就在眼前：过去一两百年里，许多国家的平均身高增长了十几厘米，而人群的基因库根本没来得及大变——变的是营养、卫生和童年疾病负担。遗传率描述的是“此时此刻这群人为什么彼此不同”，它完全管不了“换一个环境，大家会变成什么样”。谁用遗传率论证“改变环境没用”，谁就是把温度计当成了暖气片说明书。
\end{wuqu}

\begin{shiyan}
\textbf{家族性状小调查}（安全，只需要纸笔和家人的配合）。

材料：一张纸、一支笔，以及愿意配合你的家人（最好包括祖父母辈）。

步骤：列出 4～6 个容易观察的性状，例如：能不能卷舌、耳垂是贴着还是悬垂、大拇指最后一节能不能明显后弯、习惯用哪只手、是否喝牛奶后容易腹胀。逐个询问并记录每个家庭成员的情况，画成一张简单的家谱图（方框代表男性、圆圈代表女性，填色代表“有这个性状”）。

观察点：①有没有哪个性状呈现清晰的“父母有、孩子常有”模式？②有没有同卵或亲兄弟姐妹不一样的情况？③哪个性状你怀疑受环境影响很大（比如用手习惯可能受教育纠正）？

安全提示：无风险。但有一条科学诚信提示——查完之后你会明白，连“卷舌由一对等位基因决定”这种写进过教科书的说法，真实数据也并不支持它那么简单。你的调查不需要得出“结论”，它的目的是让你亲手摸到“基因到性状之间那段距离”。
\end{shiyan}

\section{再读那份基因报告}

现在回到开场那管唾液。你已经有了给报告逐句打分的工具。

“你的祖源成分约有多少来自东亚”——这类结论是与参考人群比对的统计结果，靠谱，但“成分”会随参考数据库更新而改动，它不是刻在石头上的血统证书。

“你喝酒容易上脸”——这背后常常是一个酶基因（乙醛脱氢酶）的真实变异，一个碱基换一个氨基酸，和镰状案例同款逻辑，预测力相当强。链短、效应大、外显率高，这是基因报告最能打的领域。

“你是爆发力型运动员”——警惕。运动能力是几百个基因加上训练、营养、伤病史、心理的合奏，单个位点的贡献微乎其微。这句话的诚实版本应该是：“你的某个位点在统计上与某类运动员略有关联，效应小到对你个人几乎没有指导意义。”

“你患某病的风险比平均高 $30\%$”——先问基线。再问这个评分是按哪个人群算的。最后记住：风险是群体的赌注，不是你个人的剧本，而生活方式、筛查和医疗，恰恰能改写剧本的上演方式。

你看，这份报告真正的价值，不是“剧透你的人生”，而是给你几张概率不同的牌。怎么打，还看你和你所处的环境。

\section{“坏基因”为什么没有被清除}

留一个让你睡不着的问号作伏笔。

镰状等位基因 $\mathrm{Hb^{S}}$ 在纯合时会致人于重病——按直觉，这样的“坏基因”应该一代比一代少，最终从人群里消失。可现实是，在非洲、地中海、印度的许多地区，它的频率高达 $10\%$ 甚至更高。为什么？

答案你已经见过一半：地图上那片与疟疾重合的区域。杂合子 $\mathrm{Hb^{A}Hb^{S}}$ 既不贫血，又对重症疟疾有抵抗力——在疟疾横行的环境里，他们比两种纯合子都活得好。一个等位基因，同时是“病”和“药”，全看它有几份、以及环境里有没有疟原虫。基因的“好”与“坏”，原来不是序列本身的属性，而是序列与环境谈判的结果。

谁在做这场谈判？谈判的规则是什么？群体中等位基因的频率为什么会一代代改变？——那是第十篇的开场戏，第 76 章，自然选择。我们很快回来。

\begin{tiaozhan}
把本章的直觉写成一行公式，就是数量遗传学的出发点。一个群体里某性状的表型差异（方差 $V_P$）可以拆成几份：
\[V_P = V_G + V_E + V_{G\times E}\]
其中 $V_G$ 来自基因差异，$V_E$ 来自环境差异，$V_{G\times E}$ 是“基因与环境互相搭配”的部分——比如同样的营养条件下，某个基因版本的孩子长得更快。广义遗传率定义为 $V_G$ 占 $V_P$ 的比例。这行式子立刻说明三件事：其一，遗传率是\textbf{比值}，分子里分母里都写着“这个群体、这个环境”，换群体换环境它就变；其二，$V_{G\times E}$ 的存在意味着“基因决定多少、环境决定多少”有时根本不能分账，因为账是合着记的；其三，即使 $V_{G\times E}$ 为零、遗传率等于 $1$，也不表示环境无能为力——那两块田就是证明。数学没有让结论变复杂，它只是把你已经理解的道理钉死，让人无法偷换概念。跳过本框不影响主线。
\end{tiaozhan}

\section*{练一练}

\begin{sikaoti}
\item 画一张信息流图：从 $\mathrm{Hb^{S}}$ 等位基因出发，经过 mRNA、蛋白质、红细胞、血管，到患者感受到的疼痛。在每条箭头旁注明：这一环由什么物理或化学规则驱动？哪一环最容易被环境（如饮水、感染、医疗）改变？
\item 一位朋友拿着基因报告说：“我的身高遗传率是 $80\%$，所以我长到一米七五，其中一米四归基因，剩下的归环境。”指出这句话里的两处错误，并用“两块田”的思想实验向他解释正确读法。
\item 纯合 $\mathrm{Hb^{S}Hb^{S}}$ 的外显率接近 $100\%$，但患者症状轻重差异很大。用“表现度”和本章的链条图解释：同样的基因型，为什么演出不同的剧情？请至少举出链条上两个可能产生差异的环节。
\item 教科书长期说“卷舌由一对等位基因决定，会卷是显性”。你的家族调查或班里统计如果发现了反例（比如父母都会卷、孩子不会），有哪几种解释？哪些解释与环境有关、哪些与基因互作有关？
\item 多基因风险评分在人群 A 上训练得出，用于人群 B 时准确性下降。用本章的“基因—环境—网络”模型解释为什么会这样。（提示：两个人群的等位基因频率、连锁结构、环境分布都可能不同。）
\item 某地区 $\mathrm{Hb^{S}}$ 的频率是 $10\%$，当地疟疾横行；另一个无疟疾地区该频率接近零。先不查资料，只用本章最后一节的逻辑预测：如果疟疾被彻底消灭，几代人之后 $\mathrm{Hb^{S}}$ 的频率会往哪个方向变？为什么？（第 76 章会给出完整工具。）
\end{sikaoti}
